\documentclass[11pt]{book}
\usepackage{ dblfloatfix,enumitem,graphicx,adjustbox,fancyhdr,color,multicol,bm,amsthm,amsmath}
\usepackage[hidelinks]{hyperref}
\usepackage{geometry}
\geometry{	%--------------تعیین حاشیه صفحه--------------
	a4paper,
	top=1.75cm,
	bottom=1cm,
	inner=1.25cm,
	outer=0.75cm	
}

\makeatletter
\newcommand{\mytitle}{\@title}
\makeatother
\pagestyle{fancy}

% \renewcommand{\chaptermark}[1]{\markboth{#1}{}}
\renewcommand{\sectionmark}[1]{\markboth{#1}{}}
\chead{}
\fancyhf{}
\fancyhead[CO,CE]{معادله درجه دوم}
\fancyhead[LO,RE]{\thepage}

\renewcommand{\adjboxvtop}{8mm}
%--------------تعیین فاصله بین تست ها--------------
\newcommand{\sol}{\bigskip \bigskip \bigskip \bigskip \bigskip \bigskip }

\newcommand{\riazi}[1]{\textbf{$\bm{#1}$}}
\newcommand{\quenstion}[1]{\item \adjustbox{valign=t}{\includegraphics[width=\linewidth]{../Image/#1}}
	
	\sol
}

\newcommand{\quz}[1]{\item \adjustbox{valign=t}{\includegraphics[width=\linewidth]{../Image/#1}}}

\usepackage{xepersian}	%تنظیمات فارسی نویسی
\settextfont{XB Niloofar}	%می باشد font فونت ها در داخل پوشه 
\setdigitfont{Yas}


\title{ جزوه ریاضیات و حسابان}
\author{ابوالفضل ا. | علیرضا ک.}

% -------------------------------------------------------------

%\newenvironment{mcquestions}{}{}
\newcommand{\test}[1]{\textbf{\item #1}}
\newsavebox{\answera}
\newsavebox{\answerb}
\newsavebox{\answerc}
\newsavebox{\answerd}
\newcommand{\answer}[4]{%
	\sbox\answera{ #1}%
	\sbox\answerb{ #2}%
	\sbox\answerc{ #3}%
	\sbox\answerd{ #4\sol}%
	\ifdim\wd\answera<0.165\textwidth%
	\ifdim\wd\answerb<0.165\textwidth%
	\ifdim\wd\answerc<0.165\textwidth%
	\ifdim\wd\answerd<0.165\textwidth%
	\begin{multicols}{4}
		\begin{itemize}
			\item[1)] \usebox{\answera}
			\item[2)] \usebox{\answerb}
			\item[3)] \usebox{\answerc}
			\item[4)] \usebox{\answerd}
		\end{itemize}
	\end{multicols}
	\else%
	\ifdim\wd\answera<0.4\textwidth%
	\ifdim\wd\answerb<0.4\textwidth%
	\ifdim\wd\answerc<0.4\textwidth%
	\ifdim\wd\answerd<0.4\textwidth%
	\begin{multicols}{2}
		\begin{itemize}
			\item[1)] \usebox{\answera}
			\item[3)] \usebox{\answerc}
			\item[2)] \usebox{\answerb}
			\item[4)] \usebox{\answerd}
		\end{itemize}
	\end{multicols}
	\else
	\begin{itemize} 
		\item[1)]  #1
		\item[2)]  #2
		\item[3)]  #3
		\item[4)]  #4
	\end{itemize}
	\fi
	\else
	\begin{itemize} 
		\item[1)] #1
		\item[2)] #2
		\item[3)] #3
		\item[4)] #4
	\end{itemize}
	\fi
	\else
	\begin{itemize} 
		\item[1)] #1
		\item[3)] #2
		\item[2)] #3
		\item[4)] #4
	\end{itemize}
	\fi
	\else
	\begin{itemize} 
		\item[1)] #1
		\item[3)] #2
		\item[2)] #3
		\item[4)] #4
	\end{itemize}
	\fi\fi
	\else
	\ifdim\wd\answera<0.4\textwidth%
	\ifdim\wd\answerb<0.4\textwidth%
	\ifdim\wd\answerc<0.4\textwidth%
	\ifdim\wd\answerd<0.4\textwidth%
	\begin{multicols}{2}
		\begin{itemize}
			\item[1)] \usebox{\answera}
			\item[3)] \usebox{\answerc}
			\item[2)] \usebox{\answerb}
			\item[4)] \usebox{\answerd}
		\end{itemize}
	\end{multicols}
	\else
	\begin{itemize} 
		\item[1)] #1
		\item[3)] #2
		\item[2)] #3
		\item[4)] #4
	\end{itemize}
	\fi
	\else
	\begin{itemize} 
		\item[1)] #1
		\item[3)] #2
		\item[2)] #3
		\item[4)] #4
	\end{itemize}
	\fi
	\else
	\begin{itemize} 
		\item[1)] #1
		\item[3)] #2
		\item[2)] #3
		\item[4)] #4
	\end{itemize}
	\fi
	\else
	\begin{itemize} 
		\item[1)] #1
		\item[3)] #2
		\item[2)] #3
		\item[4)] #4
	\end{itemize}
	\fi\fi
	\else
	\ifdim\wd\answera<0.4\textwidth%
	\ifdim\wd\answerb<0.4\textwidth%
	\ifdim\wd\answerc<0.4\textwidth%
	\ifdim\wd\answerd<0.4\textwidth%
	\begin{multicols}{2}
		\begin{itemize}
			\item[1)] \usebox{\answera}
			\item[3)] \usebox{\answerc}
			\item[2)] \usebox{\answerb}
			\item[4)] \usebox{\answerd}
		\end{itemize}
	\end{multicols}
	\else
	\begin{itemize} 
		\item[1)] #1
		\item[3)] #2
		\item[2)] #3
		\item[4)] #4
	\end{itemize}
	\fi
	\else
	\begin{itemize} 
		\item[1)] #1
		\item[3)] #2
		\item[2)] #3
		\item[4)] #4
	\end{itemize}
	\fi
	\else
	\begin{itemize} 
		\item[1)] #1
		\item[3)] #2
		\item[2)] #3
		\item[4)] #4
	\end{itemize}
	\fi
	\else
	\begin{itemize} 
		\item[1)] #1
		\item[3)] #2
		\item[2)] #3
		\item[4)] #4
	\end{itemize}
	\fi\fi
	\else
	\ifdim\wd\answera<0.4\textwidth%
	\ifdim\wd\answerb<0.4\textwidth%
	\ifdim\wd\answerc<0.4\textwidth%
	\ifdim\wd\answerd<0.4\textwidth%
	\begin{multicols}{2}
		\begin{itemize}
			\item[1)] \usebox{\answera}
			\item[3)] \usebox{\answerc}
			\item[2)] \usebox{\answerb}
			\item[4)] \usebox{\answerd}
		\end{itemize}
	\end{multicols}
	\else
	\begin{itemize} 
		\item[1)] #1
		\item[3)] #2
		\item[2)] #3
		\item[4)] #4
	\end{itemize}
	\fi
	\else
	\begin{itemize} 
		\item[1)] #1
		\item[3)] #2
		\item[2)] #3
		\item[4)] #4
	\end{itemize}
	\fi
	\else
	\begin{itemize} 
		\item[1)] #1
		\item[3)] #2
		\item[2)] #3
		\item[4)] #4
	\end{itemize}
	\fi
	\else
	\begin{itemize} 
		\item[1)] #1
		\item[3)] #2
		\item[2)] #3
		\item[4)] #4
	\end{itemize}
	\fi\fi
\sol}
% -------------------------------------------------------------

\begin{document}
\maketitle

\thispagestyle{empty}

\bigskip



\begin{center}
	\begin{minipage}{17cm}
{\large 
	\setlength{\parindent}{1.5em}
	
	\textbf{{\huge مقدمه }}\bigskip
	
	\indent
این مجموعه، شامل تست‌هایی از مباحث معادلات درجه دوم، گنگ و گویا و همچنین کمی سهمی است، عموم سوالات در خصوص معادله درجه دوم است، توجه ویژه‌ای بر روابط بین ریشه‌های معادلات درجه دوم شده، سوالاتی درخصوص جمع و ضرب ریشه‌ها و...
\smallskip

برای جمع‌آوری این سوالات سعی کردیم از منابع متنوع و متعددی کمک بگیریم. سوالات عموماً از آزمون‌های آزمایشی کنکور سراسری (قلمچی، گزینه دو، گاج، ماز و...)، آزمون‌های مدارس برتر و کنکور سراسری دو-سه دهه اخیر و همچنین المپیادهای ریاضی انتخاب می‌شوند.\smallskip

این مجموعه، هنوز تکمیل نشده و ما در تلاش مستمر برای تکمیل و بهبود و جامع‌تر کردن‌ آن هستیم، شما می‌توانید آخرین نسخه این مجموعه را با اسکن کد زیر دریافت کنید.

\bigskip

\begin{latin}   


This pamphlet includes multiple-choice questions from Quadratic Equations, Irrational Equations, Rational Equations and Parabola, most of the questions are about Quadratic Equations, Special attention has been paid to the relationship between the roots of quadratic equations ($\alpha$,$\beta$). questions about sum, product, etc of roots.



Various and numerous resources have been used to compile this pamphlet, The sources are mostly from test of institutions such as Maze, Qalamchi, Gaj, Qozine 2, etc, test of high schools and the national entrance exam (konkur) of the last two decades and also Mathematics Olympiad.%}

%\lr{
The pamphlet hasn't been completed yet and we are constantly striving to complete, improve and make it more comprehensive.
%}

\hfill - October, 2022 - 
\end{latin}

}\end{minipage}
\end{center}


\sol 

\begin{figure*}[bp]
\begin{center}
	\includegraphics[width=3cm]{../Image/QrCode/github-page}
\end{center}

\bigskip

\hfill 
\end{figure*}

\newpage

\clearpage
\pagenumbering{arabic}
	
\begin{enumerate}[label={\textbf{\arabic*-}}]
	\quenstion{1}
	\quenstion{2}
	\quenstion{3}
	\quenstion{4}
	\quenstion{5}
	\quenstion{6}
	\quz{7} \bigskip \bigskip
	\quenstion{8}
	\quenstion{9}
	
	\quenstion{10}
	\quenstion{11}
	\quenstion{12}
	\quenstion{13}
	\quenstion{14}
	\quenstion{15}
	\quenstion{16}
	\quenstion{17}
	\quenstion{18}
	\quenstion{19}
	
	\quenstion{20}
	\quenstion{21}
	\quenstion{22}
	\quenstion{23}
	\quenstion{24}
	\quenstion{25}
	\quenstion{26}
	\quz{27} \bigskip \bigskip
	\quenstion{28}
	\quenstion{29}
	
	\quenstion{30}
	\quenstion{31}
	\quenstion{32}
	\quenstion{33}
	\quenstion{34}
	\quenstion{35}
	\quz{36} \bigskip \bigskip
	\quenstion{37}
	\quenstion{38}
	\quz{39} \bigskip \bigskip
	
	\quenstion{40}
	\quenstion{41}
	\quenstion{42}
	\quenstion{43}
	\quz{44} \bigskip \bigskip
	\quz{45} \bigskip \bigskip \bigskip
	\quenstion{46}
	\quenstion{47}
	\quenstion{48}
	\quenstion{49}
	
	\quenstion{50}
	\quenstion{51}
	\quenstion{52}
	\quenstion{53}
	\quenstion{54}
	\quenstion{55}
	\quenstion{56}
	\quenstion{57}
	\quz{58} \bigskip \bigskip
	\quenstion{59}
	
	\quenstion{60}
	\quenstion{61}
	\quenstion{62}
	\quenstion{63}
	\quenstion{64}
	\quenstion{65}
	\quenstion{66}
	\quenstion{67}
	\quenstion{68}
	\quenstion{69}
	
	\quz{70} \bigskip \bigskip
	\quenstion{71}
	\quenstion{72}
	\quenstion{73}
	\quenstion{74}
	\quenstion{75}
	\quenstion{76}
	\quenstion{77}
	\quenstion{78}
	\quenstion{79}
	
	\quenstion{80}
	\quz{81} \bigskip \bigskip
	\quenstion{82}
	\quz{83} \bigskip \bigskip \bigskip
	\quenstion{84}
	\quenstion{85}
	\quenstion{86}
	\quenstion{87}
	\quenstion{88}
	\quenstion{89}
	
	\quenstion{90}
	\quz{91} \bigskip \bigskip
	\quenstion{92}
	\quenstion{93}
	\quz{94} \bigskip \bigskip
	\quenstion{95}
	\quenstion{96}
	\quenstion{97}
	\quenstion{98}
	\quenstion{99}
	
	\quenstion{100}
	\quenstion{101}
	\quenstion{102}
	\quz{103} \bigskip \bigskip
	\quenstion{104}
	\quenstion{105}
	\quz{106} \bigskip \bigskip
	\quenstion{107}
	\quenstion{108}
	\quz{109} \bigskip \bigskip
	
	\quz{110} \bigskip
	\quenstion{111}
	\quenstion{112}
	\quenstion{113}
	\quenstion{114}
	\quenstion{115}
	\quenstion{116}
	\quenstion{117}
	\quenstion{118}
	\quenstion{119}
	
	\quenstion{120}
	\quenstion{121}
	\quz{122} \bigskip \bigskip
	\quenstion{123}
	\quz{124} \bigskip \bigskip
	\quenstion{125}
	\quenstion{126}
	\quenstion{127}
	\quenstion{128}
	\quenstion{129}
	
	\quenstion{130}
	\quenstion{131}
	\quenstion{132}
	\quenstion{133}
	\quenstion{134}
	\quenstion{135}
	\quz{136} \bigskip \bigskip \bigskip
	\quz{137} \bigskip \bigskip
	\quenstion{138}
	\quz{139} \bigskip \bigskip \bigskip
	
	\quenstion{140}
	\quenstion{141}
	\quenstion{142}
	\quz{143} \bigskip \bigskip
	\quz{144} \bigskip \bigskip \bigskip
	\quenstion{145}
	\quenstion{146}
	\quenstion{147}
	\quenstion{148}
	\quenstion{149}
	
	\quenstion{150}
	\quenstion{151}
	\quenstion{152}
	\quenstion{153}
	\quenstion{154}
	\quz{155} \bigskip \bigskip
	\quenstion{156}
	\quenstion{157}
	\quenstion{158}
	\quenstion{159}
	
	\quenstion{160}
	\quenstion{161}
	\quenstion{162}
	\quenstion{163}
	\quenstion{164}
	\quenstion{165}
	\quenstion{166}
	\quenstion{167}
	\quenstion{168}
	\quenstion{169}
	
	\quenstion{170}
	\quenstion{171}
	\quenstion{172}
	\quz{173} \bigskip \bigskip
	\quz{174} \bigskip \bigskip \bigskip
	\quenstion{175}
	\quenstion{176}
	\quenstion{177}
	\quenstion{178}
	\quz{179} \bigskip \bigskip \bigskip \bigskip
	
	\quz{180} \bigskip \bigskip \bigskip \bigskip
	\quz{181} \bigskip \smallskip
	\quz{182} \bigskip \smallskip
	\quenstion{183}
	\quenstion{184}
	\quenstion{185}
	\quenstion{186}
	\quenstion{187}
	\quenstion{188}
	\quz{189} \bigskip \bigskip
	
	\quenstion{190}
	\quenstion{191}
	\quenstion{192}
	\quenstion{193}
	\quz{194} \bigskip \bigskip
	\quenstion{195}
	\quenstion{196}
	\quenstion{197}
	\quenstion{198}
	\quz{199} \bigskip \bigskip \bigskip
		
	\quenstion{200}
	\quenstion{201}
	\quenstion{202}
	\quenstion{203}
	\quenstion{204}
	\quenstion{205}
	\quenstion{206}
	\quenstion{207}
	\quenstion{208}
	\quenstion{209}\newpage

\test{
به ازای چند عدد صحیح \riazi{m} نمودار تابع 
\riazi{f(x)=(2m^2+6m)x^2 -2mx -1}
از ناحیه اول \underline{نمی‌گذرد}؟
}
\answer{صفر}{۱}{۲}{۳}
\test{
در معادله درجه دوم 
\riazi{x^2 - (a-2)x + (a+1)}،
کم‌ترین مقدار مجموع مربعات ریشه‌ها کدام است؟
}
\answer{۱}{۳}{۴}{۵}
	\quz{212} \bigskip \bigskip
\test{
معادله 
\riazi{x^4 - (m^2-1)x^2+3-4m}،
چهار ریشهٔ حقیقی دارد که مجموع مربعات آن‌‌ها برابر ۳۰ است.چند مقدار برای \riazi{m} وجود دارد؟
}
\answer{صفر}{40}{۲}{۴}
	\quenstion{214} \newpage
\test{
اگر 
\riazi{\alpha} و \riazi{\beta}
ریشه‌های معادلهٔ 
\riazi{x^2 +3x -1 =0}
باشند، کدام معادلهٔ زیر ریشه‌هایش 
\riazi{\beta^2 + 3\beta^2+1}
و
\riazi{\alpha^2 + 3\alpha^2+1}
می‌باشند؟
}
\answer{$x^2+x-3=0$}{$2x^2-x+3=0$}{$x^2-3x+1=0$}{$3x^2-x-1=0$}
	\quenstion{216}
	\quenstion{217}
	\quz{218} \bigskip \bigskip \bigskip
	\quenstion{219}
	\quenstion{220}
	\quenstion{221}
	\quenstion{222}
	\quenstion{223}

	\test{
اگر 
\riazi{\alpha}
 و 
\riazi{\beta} 
اعداد طبیعی و ریشه‌های معادله 
\riazi{x^2 -(\alpha^2 + \beta^2 -12)x+\beta+\alpha-1=0}
 باشد، مقدار 
\riazi{\beta+\alpha}
  کدام است؟
	}
	\answer{2}{5}{9}{12}
	\newpage
	\test{
	به ازای دو مقدار \riazi{a}، یک ریشه معادله 
	\riazi{3x^2 -ax +4 =0}،
	سه برابر دیگر است، اختلاف این دو مقدار \riazi{a} کدام است؟
}\answer{8}{9}{16}{18}

	
	\test{
یک تاس را پرتاب می‌کنیم و عدد رو شده را به جای \riazi{m}	در معادله
\riazi{x^2 - mx +(m-1)=0}
قرار می‌دهیم. احتمال آن‌که معادله دارای دو ریشه حقیقی مثبت و متمایز باشد، کدام است؟
}\answer{$\dfrac16$}{$\dfrac13$}{$\dfrac12$}{$\dfrac23$}
	
	\quenstion{223}
	\quenstion{224}
	\quenstion{225}
	\quenstion{226}
	\quenstion{227}
	\quenstion{228}
	\quenstion{229}
	\quenstion{230}
	\quenstion{231}
	
	\test{
	اگر مجموعه جواب نامعادلهٔ
	\riazi{\dfrac{((m^2-1)x^2-4mx+4)(x-3\sqrt{x}+2)}{2x-3}\le 0}
،
به ازای 
\riazi{x>\dfrac{3}{2}}،
بازهٔ 
\riazi{[2,4]}
 باشد؛ مقدار \riazi{m}، کدام است؟
}\answer{-2}{صفر}{۱}{۲}


\quenstion{232}
\quz{233}\answer{$1<m<5$}{$2<m<5$}{$2<m<4$}{$2<m<6$}
\quz{234}\answer{$m<-6$}{$m>3$}{$0<m<3$}{$3<m<6$}\newpage
\quz{235}\answer{$-1<m<0$}{$m<0$}{$2<m<8$}{$m>8$}
\quz{236}\answer{$x^2+3x+1=0$}{$x^2-3x+1=0$}{$x^2+5x+2=0$}{$x^2-5x+2=0$}
\quz{237}\answer{$a<-4$}{$a>-4$}{$a<4$}{$a>4$}
\quz{238}\answer{$-\dfrac{9}{5}$}{$1$}{$1,-\dfrac{9}{5}$}{$-1,\dfrac{9}{5}$}
\quz{239}\answer{$-4$}{$4,-12$}{$-4,12$}{$12$}\newpage

\quz{240}\answer{$-4$}{$-2$}{$2$}{$4$}
\quz{241}\answer{$-2$}{$-1$}{$1$}{$2$}
\quz{242}\answer{$-1<m<5$}{$-1<m<4$}{$-2<m<4$}{$-1<m<5$}
\quz{243}\answer{$-\dfrac{3}{2}<m<5$}{$0<m<2$}{$\dfrac{3}{2}<m<\dfrac{5}{2}$}{$\dfrac{3}{2}<m<2$}
\quz{244}\answer{$m\ge1$}{$m<2$}{$1\le m<2$}{هیچ مقدار $m$}\newpage
\quz{245}\answer{$-2$}{$-1$}{$\dfrac{2}{3}$}{$\dfrac{4}{2}$}
\quz{246}\answer{$-14$}{$-12$}{$-8$}{$-6$}
\quenstion{247}
\quenstion{248}
\quenstion{249}

\quenstion{250}
\quenstion{251}
\quenstion{252}
\quenstion{253}
\quenstion{254}
\quenstion{255}
\quenstion{256}
\quenstion{257}
\quenstion{258}
\quenstion{259}

\quenstion{260}
\quenstion{261}
\quenstion{262}
\quenstion{263}
\quenstion{264}
\quenstion{265}
\quenstion{266}
\quenstion{267}
\quenstion{268}
\quenstion{269}

\quenstion{270}
\quenstion{271}
\quenstion{272}
\quenstion{273}
\quenstion{274}
\quenstion{275}
\quenstion{276}
\quenstion{277}
\quz{278} \bigskip \bigskip
\quenstion{279}

\quenstion{280}
\quenstion{281}
\quenstion{282}
\quenstion{283}
\quenstion{284}
\quenstion{285}
\quenstion{286}
\quenstion{287}
\quenstion{288}
\quenstion{289}

\quenstion{290}
\quenstion{291}
\quenstion{292}\newpage
\quz{293}\answer{$m<0$}{$m>4$}{$-1<m<4$}{$-2<m<6$}
\quz{294}\answer{$-1$}{2}{4}{6}
\quz{295}\answer{$3.5$}{4}{$4.5$}{5}
\quz{296}\answer{2}{3}{4}{6}
\quz{297}\answer{$(-\infty,2) \cup (6,+\infty)$}{$(-\infty,3) \cup (4,+\infty)$}{$(2,6)$}{$(3,4)$}\newpage
\quz{298}\answer{1}{2}{3}{4}
\quz{299}\answer{$-1$}{$-10$}{8}{6}

\quz{300}\answer{$-6,-2$}{$-6,2$}{$6,-1$}{$-6,1$}
\quz{301}\answer{$-4$}{$-3$}{3}{4}
\quz{303}\answer{$m<-4$}{$m>4$}{$-4<m<4$}{$4<m<9$}

\newpage

\test{
معادلهٔ
\riazi{\dfrac{1}{\sqrt{2-x}+2}- \dfrac{1}{2- \sqrt{2-x}} = \dfrac{2-x}{5\sqrt{2-x}}}
چند ریشه دارد؟
}\answer{صفر}{1}{2}{3}

\test{
اگر سهمی
\riazi{y=-ax^2+ax+2}
روی سهمی
\riazi{y=2bx^2-bx-1}
قرار داشته باشد و برعکس، مقدار
\riazi{a-b}
چقدر است؟
}\answer{$-6$}{$6$}{$-18$}{18}

\test{
اگر
\riazi{\alpha} و \riazi{\beta}
ریشه‌های معادلهٔ 
\riazi{4x^3+kx^2-9x-2=0}، \riazi{\alpha + \beta = 1} و \riazi{\alpha\beta = -2}
باشد، مقدار 
\riazi{k} چقدر است؟
}\answer{$-\dfrac{27}{5}$}{$\dfrac{27}{5}$}{$-3$}{3}


\test{
فاصله نقطهٔ تلاقی منحنی‌های
\riazi{2y=x^2} و \riazi{x=\sqrt{y+3}-\sqrt{y-3}}
با مبدأ مختصات، کدام است؟
}\answer{$\sqrt{3}$}{$\sqrt{6}$}{$2\sqrt{3}$}{$\sqrt{15}$}

\test{
معادلهٔ
\riazi{\dfrac{\sqrt{x+1}}{\sqrt{x-1}+3} - \dfrac{\sqrt{x+1}}{3-\sqrt{x-1}} = \dfrac{x-1}{\sqrt{x-1}}}
چند ریشهٔ مثبت دارد؟
}\answer{صفر}{1}{2}{3}
\newpage
\test{
به ازای چند مقدار
\riazi{y=ax^2+(3+2a)x}
از ناحیه سوم محور‌های مختصاتی \underline{نمی‌گذرد}؟
}\answer{هیچ مقدار $a$}{تمام مقادیر $a$}{1}{2}


\quenstion{304}
\quenstion{305}
\quenstion{306}
\quenstion{307}
\quenstion{308}
\quenstion{309}
\end{enumerate}

\end{document}